% !PRINTED 497
\weberpage{497}

The product of these homogeneous functions cannot, by what was proved above, vanish if none of the factors vanishes, and its degree is equal to the sum of the degrees of the individual factors.

The numbers of the field $\Omega$ are included among the functions. One obtains them by letting either the degree or the number of variables sink down to zero.

Concerning the integral functions in $\Omega$ we have derived several important theorems in the first volume, of which we here single out the theorem \S.\,141, IV:

1. \gesperrt{Among the integral functions in $\Omega$ a distinction must be made between reducible and irreducible ones, of which the former can be represented as a product of several integral functions in $\Omega$, the others not. The reducible functions $\varphi$ can be decomposed into a finite number of irreducible factors, which are themselves integral functions in $\Omega$, and the irreducible factors of $\varphi$ are, apart from constant factors, i.\,e. numbers in $\Omega$, uniquely determined by $\varphi$ itself.}

As special cases, among the integral functions in $\Omega$ are also contained the integral functions in the field of the rational numbers $R$:
\weq{2}{\Phi(x,y,z,\ldots) = \sum a\,x^r y^s z^t \ldots,}
in which the coefficients $a$ are rational numbers.

If these coefficients are \gesperrt{integers without a common divisor}, then this function is called an \gesperrt{original or primitive} one, of which we have proved the following theorem in Bd. I, \S.\,2:

2. \gesperrt{The product of two primitive functions is again a primitive function.}

If the coefficients of the function (2), which we shall for the moment denote by
\[ a_0, a_1, a_2, \ldots \]
are integers with the greatest common divisor $m$, then, if $m>1$, $\Phi$ is an imprimitive \gesperrt{integral integer-coefficient function of divisor $m$, % [sic] "ganze ganzzahlige" so printed in the German
and the divisor of a primitive function is $= 1$.}
